\chapter{45}\label{section-44}

Der diensthabende Arzt auf der Notfallstation des \emph{Ospedale} kam nach einer guten Dreiviertelstunde den Flur entlang auf den Commissario zu. Mantovani erwartete die ersten Befunde. Dabei hatte er ununterbrochen sein Handy am Ohr.

«Nein, leider nicht. Kein Handy, keine Tasche, nichts, wir haben noch keine Personalien.»

Mantovani sprach gereizt. Er hatte kaum geschlafen. Bereits vergangene Nacht war er in den ersten Morgenstunden geweckt worden. Auch ein Notfall.

«Ja, exakt, alle Vermisstenmeldungen,» seine Stimme bekam einen fordernden Unterton, «Rosario bleibt neben dem Telefon sitzen, bis ich zurück bin.»

Aus einem Zimmer neben Mantovani drang ein Schrei in den Flur.

«Und, Melik, suche bitte den Platz oben am Felsen akribisch nach möglichen Spuren ab.»

Mantovani hörte kurz zu.

«Hm, ja, genau da. Nimm Rossi mit.»

Als der Arzt vor ihm stand, verabschiedete er sich.

«Ciao, Melik, danke, ich bin gegen 10:00 Uhr in der Quästur, bis dann.»

«Es sieht klar nach Sturz aus,» begann der Arzt. Er sah jung aus, «sie
muss aus beachtlicher Höhe runtergefallen sein, auf einen Baum, sagten
Sie?»

Mantovani nickte und der Arzt fuhr engagiert fort.

«Verletzung am Hinterkopf, eine offene Wunde, keine Schädelfraktur auf
den ersten Blick, Blutverlust, aber nicht lebensgefährlich, Kratzspuren,
als einzelne Risse am Handballen rechts plus einen neun Millimeter
langen, dünnen Holzsplitter mittendrin, Durchmesser ca. zwei Millimeter,
ferner eine schwere Gehirnerschütterung mit Bewusstlosigkeit,
gebrochener Arm rechts, Rückenverletzung wohl durch Aufprall. Das muss
ich genauer untersuchen, die Wirbelsäule könnte verletzt sein. Wir haben
sie ins künstliche Koma gesetzt, wir brauchen Zeit. Die Herzfrequenz ist
aktuell ausgeglichen. Sie war leicht unterkühlt. Sie hat mit Sicherheit
länger da oben gelegen.»

Der Arzt machte ein fleissiges Gesicht.

Mantovani hatte zu seinen Ausführungen mehrmals genickt. Er kniff seine
Augen zusammen.

«Wie lange schätzen Sie, lag sie da?»

«Schwer zu sagen, vielleicht zwei, höchstens vier Stunden. Sie hat Glück
gehabt, so mild wie die Nächte sind. Wie gesagt, sie muss bewusstlos
dort oben gelegen haben.»

«Gut, danke \emph{Dottore}. Noch eine Frage: Haben Sie irgend etwas
entdeckt, das Hinweise auf ihre Identität geben könnte, Tattoos,
irgendwas anderes?»

Mantovani fixierte ihn. Der Arzt blickte nachdenklich an die Decke,
schüttelte erst den Kopf und schob dann nach...

«...na ja, also...auf ihrem Slip stand oben \emph{Giorgia} geschrieben.
Ich habe Namen auf Frauen-Unterwäsche so noch nie gesehen. Ich weiss
nicht, ob das `ne Marke ist?»

Der Arzt lächelte und hob seinen Blick.

«Wer weiss,» sagte Mantovani und klopfte ihm auf die Schulter, «zu
diesem Holzsplitter an ihrer Hand, könnte sich die Frau irgendwo
festgehalten haben?»

«Möglich, ja, und hätte sich dabei den Holzspiess in den Handballen
gerammt.»

«Danke, Dottore, lassen Sie mir bitte die Fotos zu den verschiedenen
Verletzungen so rasch als möglich zukommen. Und bitte, melden Sie sich
bei mir, wenn Sie mehr wissen.»

Der Arzt nickte. Mantovani streckte ihm seine Visitenkarte zu, die der
in die Tasche seines Arztkittels gleiten liess. Auch er sah geschafft
aus. Er reichte dem Commissario die Hand.

Mantovani zückte sein Handy und tippte die Informationen ein. Von der
Unterwäsche-Marke versprach er sich wenig, notierte aber auch diesen
Hinweis. Er hoffte auf eine rasche Identifikation der Frau.

Sein Tag würde lang werden. Mit Wehmut fiel ihm ein, dass heute seine
jüngste Enkelin Geburtstag hatte.
